Large Language Models (LLMs) are widely used in software engineering (SE)
research and practice, yet their non-determinism, opaque training data, and
rapidly evolving models threaten the reproducibility and replicability of
empirical studies.
We address this challenge through a collaborative effort of 22 researchers,
presenting a taxonomy of seven study types that organizes how LLMs are
used in SE research, together with eight guidelines for designing
and reporting such studies.
Each guideline distinguishes requirements (must) from recommended practices
(should) and is contextualized by the study types it applies to.
Our guidelines recommend that researchers:
(1)~declare LLM usage and role;
(2)~report model versions, configurations, and customizations;
(3)~document the tool architecture beyond the model;
(4)~disclose prompts, their development, and interaction logs;
(5)~validate LLM outputs with humans;
(6)~include an open LLM as a baseline;
(7)~use suitable baselines, benchmarks, and metrics; and
(8)~articulate limitations and mitigations.
We complement the guidelines with an applicability matrix mapping guidelines
to study types and a reporting checklist for authors and reviewers.
We maintain the study types and guidelines online as a living resource for the
community to use and shape
(\href{https://llm-guidelines.org/}{llm-guidelines.org}).
